\documentclass[10pt,a4paper]{amsart}
\usepackage[margin=19mm]{geometry}
\usepackage{amsmath,amssymb,mathtools}
\usepackage[scr=rsfs,cal=euler]{mathalfa}
\usepackage{xfrac}
\usepackage[unicode,hidelinks,bookmarks=false]{hyperref}
\newcommand{\ie}{i.e.\ }
\newcommand{\cf}{cf.\ }
\newcommand{\resp}{respectively\ }
\newcommand{\Subsection}{\subsection}
\providecommand{\nobreakdash}{\nobreak\hbox{-}}
\setlength{\emergencystretch}{2em}
\multlinegap0pt
\usepackage{amssymb}
\usepackage{mathtools}
\usepackage{enumerate}
\newtheorem{theorem}{Theorem}
\title{Two Fefferman-type constructions involving almost Grassmann structures and path geometries}
\author{Zhangwen Guo}
\date{Source: arXiv:2509.04878v2 (CC BY 4.0)}
\begin{document}
\maketitle

\noindent\emph{Source and licence.} Two Fefferman-type constructions involving almost Grassmann structures and path geometries. Version arXiv:2509.04878v2 (CC BY 4.0), distributed by arXiv under the Creative Commons Attribution 4.0 International licence, http://creativecommons.org/licenses/by/4.0/, as recorded in the arXiv licence metadata for this exact version and shown on the abstract page https://arxiv.org/abs/2509.04878v2. Full licence notice: \texttt{licenses/arxiv-2509.04878-CC-BY-4.0.txt}.\par\smallskip

\noindent\textit{Editorial scope.} Counted unit: the article's theorem theorem (source0.tex lines 388--395). The article's complete original proof of it is delivered with it (source0.tex lines 397--411, one \texttt{proof} environment of 15 lines). No mathematics is written, repaired or re-derived: the statement and the proof are the article's own lines, in the article's own notation, and no retained line is substituted or re-flowed.  Retained uncounted as the source-local context the proof needs: lines 375--377.  Nothing was omitted from any retained span, so the package carries no omission record.  No retained line contains a citation, so the wrapper carries no bibliography and no bibliography entry.  No retained line contains a figure environment, an image-inclusion command or drawing code, so no image is delivered and the package declares no asset dependency.  Editorial formatting only, in addition: an AMS \texttt{amsart} preamble that restates the article's own declarations and macros used by the retained lines, namely the environments \texttt{theorem} on separate counters, together with the standard packages those lines need.  The retained environments print in this document as Theorem 1 in order of appearance, whereas the article prints the same items with the numbering of its own full document.  The complete native source of the article as distributed in the arXiv e-print for this exact version is bundled unchanged as \texttt{sources/184512/source0.tex}.
\begin{align}\label{tilde kappa}
    \tilde\kappa(j(u))=i'\circ\kappa(u)\circ\wedge^2\, \pi
\end{align}

\begin{theorem}\label[theorem]{Hol}
Notation as above. $(\tilde{\mathcal G}\to\tilde M,\tilde\omega)$ is locally the result of the Fefferman-type functor applied to a Cartan geometry of type $(G,P)$ if and only if there exists a parallel section
\[\sigma\in C^\infty(\tilde{\mathcal G},\tilde P\,\tilde w_0)^{\tilde P}=\Gamma(\tilde{\mathcal O})\subseteq\Gamma(\tilde{\mathcal W})\]
of the tractor bundle $\tilde{\mathcal W}$ such that whenever $\sigma(u)=\tilde w_0$ where $u\in\tilde{\mathcal G}$, there holds
\begin{align*}
\tilde\kappa(u)(i'(\mathfrak p),\tilde{\mathfrak g})=0.   
\end{align*}
\end{theorem}

\begin{proof}
Let us identify $H\subseteq G$ with their images under $i$ and identify their Lie algebras with their images under $i'$.

Assume that $(\tilde{\mathcal G}\to\tilde M,\tilde\omega)$ is the result of the Fefferman-type functor applied to a Cartan geometry $(\mathcal G\to M,\omega)$ of type $(G,P)$. Since $H$ fixes $\tilde w_0$, the constant map $\sigma:\mathcal G\to\{\tilde w_0\}$ is $H$-equivariant. It is extended to 
\[\tilde\sigma\in C^\infty(\tilde{\mathcal G},\tilde P\,\tilde w_0)^{\tilde P}=\Gamma(\tilde{\mathcal O})\]
in a unique way. Denote by $p:\mathcal G\to \tilde M$ the natural projection. For every tangent vector in $T\mathcal G$, there holds
\begin{align}\label{tractor connection on the distinguished tractor}
    \nabla^{\tilde{\mathcal W}}_{Tp(\,\cdot\,)}\tilde\sigma=d\,\sigma(\,\cdot\,)+\omega(\,\cdot\,)\,\sigma.
\end{align}
Here, $d\,\sigma$ denotes the directional derivative of the constant map $\sigma:\mathcal G\to\{\tilde w_0\}$ and $\omega(\,\cdot\,)\,\sigma$ denotes the induced Lie algebra action of $\mathfrak g$ on $\tilde w_0$. Both of them are zero here. In particular, $\tilde\sigma$ is a parallel section of $\tilde{\mathcal W}$. Moreover, it follows from the definition of $\tilde{\mathcal W}$ that for $u\in\tilde{\mathcal G}$, there holds $\tilde\sigma(u)=\tilde w_0$ if and only if $u\in\mathcal G$. In this case, it follows from \eqref{tilde kappa} that $\tilde\kappa(u)(i'(\mathfrak p),\tilde{\mathfrak g})=0$.

Conversely, let $\tilde\sigma\in C^\infty(\tilde{\mathcal G},\tilde P\,\tilde w_0)^{\tilde P}=\Gamma(\tilde{\mathcal O})$ be as stated in the theorem and
\[\mathcal G=\{u\in\tilde{\mathcal G}:\tilde\sigma(u)=\tilde w_0\}.\]
Observe that the natural projection $\mathcal G\to\tilde M$ admits smooth local sections and that the preimage of every point in $\tilde M$ is exactly an $H$-orbit of the fiber of $\tilde{\mathcal G}$ over the same point. Thus $\mathcal G$ is a reduction of $\tilde{\mathcal G}$ to a principal $H$-bundle. It remains to show that $\omega=\tilde\omega|_{\mathcal G}$ has values in $\mathfrak g$ so that, in particular, $(\mathcal G\to\tilde M,\omega)$ is a Cartan geometry of type $(G,H)$. Indeed, $\sigma=\tilde\sigma|_{\mathcal G}$ is the $\tilde w_0$-valued constant map. Since $\mathcal G\subseteq\tilde{\mathcal G}$, \eqref{tractor connection on the distinguished tractor} holds where $\nabla^{\tilde{\mathcal W}}\tilde\sigma$ and $d\,\sigma$ are both zero. In particular, the values of $\omega$ annihilate $\tilde w_0$. Since $G=\mathrm{Stab}_{\tilde G}(\tilde W_0)$, there holds $\mathfrak g=\mathrm{Ann}_{\tilde{\mathfrak g}}(\tilde w_0)$. In particular, the images of $\omega$ lie in $\mathfrak g$. This proves our claim. Moreover, it follows from the assumption on $\tilde\kappa$ that $(\mathcal G\to\tilde M,\omega)$ is locally the correspondence space of a Cartan geometry of type $(G,P)$. That is, $(\tilde{\mathcal G}\to\tilde M,\tilde\omega)$ is locally the result of the Fefferman-type functor applied to a Cartan geometry of type $(G,P)$. This completes the proof of the theorem.
\end{proof}

\end{document}
